\section{What the efficacy evidence shows}\label{sec:evidence}
\subsection{CBT for anxiety}
The strongest case for CBT is large and should be stated plainly. In placebo-controlled trials, CBT beats placebo on target symptoms with a moderate effect ($g=0.56$ across 41 trials, $N=2{,}843$), and roughly triples the odds of responding \citep{carpenter2018cbt}. This evidence is the reason any critique has to explain why CBT works at all. The limits are also clear: response rates average about half of patients, and they vary with the criteria used \citep{loerinc2015response}. We found no evidence in this literature that speaks to the \emph{mechanism}: trials randomise to a package, not to the cognitive-restructuring component.
\citegap{effect-size trend over publication year; waitlist versus active control; risk of bias; component and dismantling trials}

\subsection{CBT for adult ADHD}
In a randomised trial of 86 medication-treated adults with persistent ADHD symptoms, CBT outperformed an attention-matched control of relaxation with educational support: 53\% versus 23\% were responders on clinician global ratings \citep{safren2010cbt}. Two features matter for our argument. Every participant was already medicated, and the programme was skills-heavy (organising and planning, managing distractibility, procrastination) with adaptive thinking as one module among several. The Cochrane review of CBT for adult ADHD concludes that the evidence is of low quality, with short-term benefit possible but unclear \citep{lopez2018cochrane}. These findings are compatible with a reading in which the active ingredient is concrete skills and habit-building, not reappraisal, but they do not isolate the component.
